%% 05-claude.tex — JRC Ispra 2026: Claude Code for empirical research
%% Scott Cunningham, European Commission, June 3--4, 2026
%% Principle: one idea per slide. No walls of text.

\documentclass[aspectratio=169,11pt]{beamer}
\usepackage{remix}
\usepackage{madrid_theme}

\usepackage{tikz}
\usetikzlibrary{positioning,arrows.meta,calc,shapes.geometric,fit,backgrounds,decorations.pathreplacing}
\usepackage{listings}
\usepackage{xcolor}
\usepackage{booktabs}
\usepackage{colortbl}

\definecolor{harvardcrimson}{HTML}{A51C30}

\lstset{
  basicstyle=\ttfamily\footnotesize,
  keywordstyle=\color{accentdark}\bfseries,
  commentstyle=\color{warmgray}\itshape,
  stringstyle=\color{ocean},
  backgroundcolor=\color{lightbg},
  frame=single, framesep=4pt,
  rulecolor=\color{warmgray!40},
  columns=fullflexible,
  showstringspaces=false,
  breaklines=true,
  numbers=none,
  xleftmargin=0.4em, xrightmargin=0.4em,
}

\newcommand{\emphbox}[1]{%
  \begin{beamercolorbox}[rounded=true,sep=7pt,wd=\linewidth]{block body}%
    \small\color{slate} #1%
  \end{beamercolorbox}}

\begin{document}

%% ---------- TITLE ----------
\codechellatitle
  {Claude Code for Empirical Research}
  {Agents, workflows, and difference-in-differences}
  {Scott Cunningham \textperiodcentered\ Baylor University}
  {European Commission \textperiodcentered\ June 3--4, 2026}

%% ====================================================================
%% PART 1 — INTRODUCTION
%% ====================================================================

\begin{frame}{Automating empirical research is no longer hypothetical}
\vspace{0.4em}
\begin{itemize}\setlength\itemsep{0.6em}
  \item \textbf{Natural sciences:} AlphaFold predicted protein structure at scale. Nobel Prize in Chemistry, 2024.
  \item \textbf{Social sciences:} Stanford's \textit{Social Catalyst Lab} is automating end-to-end empirical pipelines --- data, design, estimation, write-up.
  \item \textbf{The question is not whether.} \highlight{It's how much $H$ you intentionally keep.}
\end{itemize}
\end{frame}

%% ----------
\begin{frame}{An agent is an LLM wrapped in a loop}
\vspace{0.6em}
\begin{center}
\begin{tabular}{lp{8cm}}
  \textbf{LLM alone} & One prompt. One reply. Done. \\[0.6em]
  \textbf{Agent}     & Same LLM, but it can use tools, observe results, and decide what to do next --- repeatedly. \\
\end{tabular}
\end{center}
\vspace{0.6em}
\begin{center}
\color{warmgray}\small
Every small decision inside that loop is recorded in the session log.\\
We can read them.
\end{center}
\end{frame}

%% ----------
\begin{frame}{Chatbot vs.\ Copilot vs.\ Agent}
\vspace{0.5em}
\begin{center}
\renewcommand{\arraystretch}{1.8}
\begin{tabular}{lll}
\toprule
              & \textbf{What it does} & \textbf{Who steers} \\
\midrule
\textbf{Chatbot}  & Answers your question & You, every turn \\
\textbf{Copilot}  & Completes what you're typing & You, inline \\
\textbf{Agent}    & Plans, acts, iterates, and verifies & You set the goal \\
\bottomrule
\end{tabular}
\end{center}
\vspace{0.5em}
\begin{center}\color{warmgray}\small
Today: agents. Claude Code is an agent.
\end{center}
\end{frame}

%% ====================================================================
%% PART 2 — THE ECONOMICS OF AI
%% ====================================================================

\begin{frame}{AI flattens the isoquants of cognitive production}
\vspace{-0.1cm}
\begin{columns}[T]
\begin{column}{0.55\textwidth}
\centering
\begin{tikzpicture}[scale=0.78]
    \draw[->] (0,0) -- (6,0) node[below right] {$H$};
    \draw[->] (0,0) -- (0,5) node[above left] {$M$};
    \draw[thick,ocean] (0.8,4.5) .. controls (1.5,2.2) and (2.5,1.4) .. (5,0.8);
    \node[ocean,font=\small] at (1.7,4.3) {Pre-AI};
    \draw[thick,harvardcrimson] (0,3) -- (5,0);
    \node[harvardcrimson,font=\small] at (3.3,1.9) {Post-AI};
    \filldraw[ocean] (3.6,1.2) circle (2.5pt);
    \draw[dashed,warmgray] (3.6,0) -- (3.6,1.2);
    \draw[dashed,warmgray] (0,1.2) -- (3.6,1.2);
    \node[ocean,right,font=\small] at (3.7,1.25) {$M^*_{\text{pre}}$};
    \filldraw[harvardcrimson] (0,3) circle (2.5pt);
    \node[harvardcrimson,left,font=\small] at (-0.15,3) {$M^*_{\text{post}}$};
\end{tikzpicture}
\end{column}
\begin{column}{0.42\textwidth}
\[ Q \;=\; f(H,\;M) \]
\begin{itemize}\small
  \item $H$ = human time and attention
  \item $M$ = machine time
\end{itemize}
\vspace{0.4cm}
\highlight{The question isn't whether to use $M$.\\It's how much $H$ you keep.}
\end{column}
\end{columns}
\end{frame}

%% ---------- Shift (no label) ----------
\begin{frame}{AI shifts the production curve up}
\vspace{-0.1cm}
\centering
\begin{tikzpicture}[scale=0.72]
    \draw[->] (0,0) -- (6.3,0) node[below right] {$H$};
    \draw[->] (0,0) -- (0,6.2) node[above left] {$Q$};
    \draw[thick,ocean,domain=0.5:5,samples=100]
        plot (\x, {2.0*ln(\x+0.5)+0.2});
    \node[ocean,right,font=\scriptsize] at (5.05,3.61) {Pre-AI};
    \draw[thick,harvardcrimson,domain=0.5:5,samples=100]
        plot (\x, {2.0*ln(\x+0.5)+2.4});
    \node[harvardcrimson,right,font=\scriptsize] at (5.05,5.81) {Post-AI};
    \filldraw[ocean] (4,3.208) circle (2.5pt);
    \filldraw[forest] (4,5.408) circle (2.5pt);
    \draw[dashed,warmgray] (4,0) -- (4,5.408);
    \draw[dashed,ocean]  (4,3.208) -- (0,3.208);
    \draw[dashed,forest] (4,5.408) -- (0,5.408);
    \node[left,font=\footnotesize,ocean]  at (0,3.208) {$Q_{\text{pre}}$};
    \node[left,font=\footnotesize,forest] at (0,5.408) {$Q_{\text{post}}$};
    \node[below,font=\footnotesize] at (4,0) {$H^*$};
\end{tikzpicture}
\end{frame}

%% ---------- Safe ----------
\begin{frame}{Safe: take time back, stay above $\bar{H}$}
\vspace{-0.1cm}
\centering
\begin{tikzpicture}[scale=0.72]
    \draw[->] (0,0) -- (6.3,0) node[below right] {$H$};
    \draw[->] (0,0) -- (0,6.2) node[above left] {$Q$};
    \draw[thick,ocean,domain=0.5:5,samples=100]
        plot (\x, {2.0*ln(\x+0.5)+0.2});
    \node[ocean,right,font=\scriptsize] at (5.05,3.61) {Pre-AI};
    \draw[thick,harvardcrimson,domain=0.5:5,samples=100]
        plot (\x, {2.0*ln(\x+0.5)+2.4});
    \node[harvardcrimson,right,font=\scriptsize] at (5.05,5.81) {Post-AI};
    \fill[forest,opacity=0.22]
        (1.0,3.208) -- plot[domain=1.0:4,samples=60] (\x, {2.0*ln(\x+0.5)+2.4})
        -- (4,3.208) -- cycle;
    \filldraw[ocean] (4,3.208) circle (2.5pt);
    \draw[dashed,ocean]    (0,3.208) -- (4,3.208);
    \draw[dashed,warmgray] (4,0) -- (4,3.208);
    \draw[dashed,warmgray] (1.0,0) -- (1.0,3.208);
    \node[left,font=\footnotesize,ocean] at (0,3.208) {$Q_{\text{pre}}$};
    \node[below,font=\footnotesize] at (4,0)   {$H^*$};
    \node[below,font=\footnotesize] at (1.0,0) {$\bar{H}$};
    \filldraw[harvardcrimson] (1.0,3.208) circle (2.5pt);
    \node[forest,font=\footnotesize\bfseries] at (2.85,3.85) {Safe};
\end{tikzpicture}
\end{frame}

%% ---------- Danger ----------
\begin{frame}{Danger: $H < \bar{H}$, skills depreciate}
\vspace{-0.1cm}
\centering
\begin{tikzpicture}[scale=0.72]
    \draw[->] (0,0) -- (6.3,0) node[below right] {$H$};
    \draw[->] (0,0) -- (0,6.2) node[above left] {$Q$};
    \draw[thick,ocean,domain=0.5:5,samples=100]
        plot (\x, {2.0*ln(\x+0.5)+0.2});
    \node[ocean,right,font=\scriptsize] at (5.05,3.61) {Pre-AI};
    \draw[thick,harvardcrimson,domain=0.5:5,samples=100]
        plot (\x, {2.0*ln(\x+0.5)+2.4});
    \node[harvardcrimson,right,font=\scriptsize] at (5.05,5.81) {Post-AI};
    \fill[harvardcrimson,opacity=0.26]
        (0.5,0) -- plot[domain=0.5:1.0,samples=50] (\x, {2.0*ln(\x+0.5)+2.4})
        -- (1.0,0) -- cycle;
    \filldraw[ocean] (4,3.208) circle (2.5pt);
    \draw[dashed,ocean]    (0,3.208) -- (4,3.208);
    \draw[dashed,warmgray] (4,0) -- (4,3.208);
    \draw[dashed,warmgray] (1.0,0) -- (1.0,3.208);
    \node[left,font=\footnotesize,ocean]  at (0,3.208)  {$Q_{\text{pre}}$};
    \node[below,font=\footnotesize] at (4,0)    {$H^*$};
    \node[below,font=\footnotesize] at (1.0,-0.05) {$\bar{H}$};
    \filldraw[harvardcrimson] (0.55,2.50) circle (2.5pt);
    \draw[dashed,harvardcrimson] (0.55,0) -- (0.55,2.50);
    \draw[dashed,harvardcrimson] (0.55,2.50) -- (0,2.50);
    \node[below,font=\scriptsize,harvardcrimson] at (0.55,-0.55) {$H_d$};
    \node[left,font=\footnotesize,harvardcrimson] at (0,2.50) {$Q_d$};
    \node[harvardcrimson,font=\footnotesize\bfseries] at (0.78,0.9) {Danger};
\end{tikzpicture}
\end{frame}

%% ====================================================================
%% PART 3 — GETTING STARTED SAFELY
%% ====================================================================

\begin{frame}{Download Claude Code from the official source only}
\vspace{0.8em}
\begin{center}
\Large\texttt{claude.com/download}
\end{center}
\vspace{0.8em}
\begin{center}
\color{warmgray}
Start with the \highlight{desktop app} --- no terminal required to begin.\\[0.4em]
Available for macOS and Windows.
\end{center}
\end{frame}

%% ----------
\begin{frame}{A new phishing attack: fake install sites}
\vspace{0.6em}
\emphbox{Fake ``Claude Code install'' pages are circulating. They ask you to copy and paste a command into your terminal. That command installs malware.}
\vspace{0.6em}
\begin{itemize}\setlength\itemsep{0.5em}
  \item The threat won't come through email
  \item It looks like a legitimate software install page
  \item \highlight{Never paste terminal commands from a website you didn't navigate to yourself}
\end{itemize}
\vspace{0.4em}
\begin{center}\color{warmgray}\small
Source: TechRepublic (2026) \textperiodcentered\ \texttt{claude.com/download} is the only safe path.
\end{center}
\end{frame}

%% ====================================================================
%% PART 4 — MIXTAPETOOLS
%% ====================================================================

\begin{frame}{MixtapeTools: skills for empirical research}
\vspace{0.8em}
\begin{center}
\Large\texttt{github.com/scunning1975/MixtapeTools}
\end{center}
\vspace{0.8em}
\begin{center}
\color{warmgray}\normalsize
A library of Claude Code skills built for researchers.\\[0.6em]
\end{center}
\begin{center}
\texttt{/newproject} \quad \texttt{/split-pdf} \quad \texttt{/referee2}\\[0.4em]
\texttt{/beautiful\_deck} \quad \texttt{/blindspot} \quad \texttt{/compiledeck}
\end{center}
\end{frame}

%% ----------
\begin{frame}[fragile]{\texttt{/newproject}: scaffold a project in seconds}
\vspace{0.4em}
\begin{lstlisting}[language=]
/newproject
\end{lstlisting}
\vspace{0.4em}
Creates:
\begin{itemize}\setlength\itemsep{0.3em}
  \item Standard directory structure (\texttt{data/}, \texttt{code/}, \texttt{output/}, \texttt{docs/})
  \item \texttt{CLAUDE.md} with project context
  \item \texttt{README.md}
\end{itemize}
\vspace{0.3em}
\begin{center}\color{warmgray}\small
Give Claude access to your folder and a chapter from the Mixtape as context.
\end{center}
\end{frame}

%% ====================================================================
%% PART 5 — ZERO SHOT IN THE BACKGROUND
%% ====================================================================

\begin{frame}{Starting a zero-shot now --- in the background}
\vspace{0.8em}
\begin{center}
\large
One prompt. No instructions. Just a research question.\\[0.8em]
\color{warmgray}\normalsize
Objective: develop a DiD study using Italian or European data.\\[0.4em]
Binary treatment. Staggered adoption.\\[0.4em]
CS estimator.
\end{center}
\vspace{0.5em}
\begin{center}\color{warmgray}\small
We'll come back to this at the end of the day and see what was produced.
\end{center}
\end{frame}

%% ====================================================================
%% PART 6 — GROWING SUBMISSIONS FROM ZERO SHOT
%% ====================================================================

\begin{frame}{AER is already seeing the wave}
\vspace{0.5em}
\begin{columns}[T]
\column{0.54\textwidth}
\begin{beamercolorbox}[rounded=true,sep=8pt,wd=\linewidth]{block body}
\small\color{slate}
\textbf{Friend 1} \textperiodcentered\ 10:17 PM\\[0.4em]
``Just learn that AER is already seeing a big inflow of AI papers lol''\\[0.4em]
``Written by AI. They are very concerned.''\\[0.4em]
``No way they can handle 20--50\% more papers''\\[0.4em]
``12 editors each handling 200 papers right now''
\end{beamercolorbox}

\column{0.43\textwidth}
\begin{beamercolorbox}[rounded=true,sep=8pt,wd=\linewidth]{block body}
\small\color{slate}
\textbf{Friend 2} \textperiodcentered\ 10:18 PM\\[0.4em]
``Haha written by AI or about AI?''\\[0.6em]
``\$1,000 submission fee!''
\end{beamercolorbox}
\vspace{0.5em}
\begin{center}
\color{warmgray}\small
This conversation happened\\this week.
\end{center}
\end{columns}
\end{frame}

%% ----------
\begin{frame}{Zurich Social Catalyst Lab: papers from the workshop}
\vspace{0.8em}
\begin{center}
\large
After the Zurich workshop (May 2026):\\[0.6em]
\color{warmgray}\normalsize
Participants built working papers using Claude Code and the methods from the course.\\[0.6em]
\highlight{Zero-shot papers are not the ceiling. They are the floor.}
\end{center}
\end{frame}

%% ====================================================================
%% PART 7 — THE CS WORKFLOW (LIVE DEMO)
%% ====================================================================

\begin{frame}{We'll do this together --- live}
\vspace{0.8em}
\begin{center}
\large
One DiD paper, built in plan mode, using CS.\\[0.6em]
\color{warmgray}\normalsize
Italian or European data. Staggered adoption. We know the estimator.\\[0.6em]
It will run \highlight{all day}.\\[0.4em]
We'll come back at the end and see where it is.
\end{center}
\end{frame}

%% ----------
\begin{frame}{The workflow: nine steps in plan mode}
\vspace{0.4em}
\begin{enumerate}\setlength\itemsep{0.35em}\small
  \item \textbf{Choose covariates} via interview process
  \item \textbf{Check balance} with balance tables
  \item \textbf{Build propensity score}
  \item \textbf{Visualize rollout} with \texttt{panelview}
  \item \textbf{Rollout graphic} with cohort counts and sample shares
  \item \textbf{Outcome by cohort} over time
  \item \textbf{Choose target parameter} (weighted ATT)
  \item \textbf{Estimate} simple average + event studies with CS
  \item \textbf{Sensitivity analysis} with Rambachan--Roth
\end{enumerate}
\end{frame}

%% ----------
\begin{frame}{Verification: the human stays in the loop}
\vspace{0.8em}
\begin{center}
\begin{tabular}{ll}
  \texttt{/referee2}        & Results in a second language / framework \\[0.5em]
  GTD harness               & Organize and track multi-step pipelines \\[0.5em]
  \texttt{/beautiful\_deck} & Compile results into a presentation \\
\end{tabular}
\end{center}
\vspace{0.6em}
\begin{center}\color{warmgray}\small
The agent does the work. You verify the design and the findings.
\end{center}
\end{frame}

%% ====================================================================
%% CLOSING
%% ====================================================================

\begin{frame}[plain]
\begin{tikzpicture}[remember picture,overlay]
\fill[cream] (current page.south west) rectangle (current page.north east);
\node[anchor=center, text width=0.84\paperwidth, align=center, text=charcoal]
  at (current page.center) {%
  \fontsize{18pt}{26pt}\selectfont
  The agent does the work.\\[0.5em]
  {\color{accent}\bfseries You set the question.}\\[0.5em]
  {\color{accent}\bfseries You verify the design.}\\[0.5em]
  {\color{accent}\bfseries You own the answer.}
};
\draw[accent, line width=1pt]
  ([yshift=1.4cm,xshift=0.25\paperwidth]current page.south west)
  -- ([yshift=1.4cm,xshift=-0.25\paperwidth]current page.south east);
\end{tikzpicture}
\end{frame}

\end{document}
